% !PRINTED 80
\weberpage{80}
%<<CONT

stimmt, abgesehen vom Vorzeichen, mit der Unterdeterminante \((n-\nu)\)ter Ordnung
\[
        A_{\alpha_{\nu+1}\ldots\alpha_n}^{\beta_{\nu+1}\ldots\beta_n}
\]
überein. Der Complex der Glieder (6) wird also bezeichnet mit
\begin{equation}
        A_{\alpha_1,\alpha_2,\ldots,\alpha_\nu}^{\beta_1,\beta_2,\ldots,\beta_\nu}
        B_{\alpha_1,\alpha_2,\ldots,\alpha_\nu}^{\beta_1,\beta_2,\ldots,\beta_\nu}.
\tag{7}
\end{equation}

Wählen wir nun für \(\alpha_1,\alpha_2\ldots\alpha_\nu\) jede Combination von \(\nu\) der Ziffern \(1,2\ldots n\), deren Anzahl (nach \S.\,7) \(B_\nu^{(n)}\) ist, so erhalten wir, indem wir \(\beta_1,\beta_2\ldots\beta_\nu\) festhalten, ebenso viele Complexe der Form (7), und jedes Glied der Determinante \(A\) kommt in einem und nur in einem dieser Complexe vor.

\gesperrt{V. Demnach erhalten wir, wenn wir alle Ausdrücke (7) summiren, die Determinante \(A\):}
\begin{equation}
        A=\sum^{\alpha}
        A_{\alpha_1,\alpha_2,\ldots,\alpha_\nu}^{\beta_1,\beta_2,\ldots,\beta_\nu}
        B_{\alpha_1,\alpha_2,\ldots,\alpha_\nu}^{\beta_1,\beta_2,\ldots,\beta_\nu}.
\tag{8}
\end{equation}

Selbstverständlich kann man auch die Combination der \(\alpha\) festhalten und in Bezug auf die \(\beta\) summiren.

Dies ist der Satz von Laplace.

Noch eine andere Darstellung der Determinante \(A\) durch die ersten und zweiten Unterdeterminanten erhält man auf folgende Weise.

Man wähle in \(A\) irgend zwei Reihen aus, die sich in einem Element, etwa in \(a_\nu^{(\mu)}\), schneiden. In jedem Gliede von \(A\) kommt ein Element mit dem unteren Index \(\nu\) und ein Element mit dem oberen Index \(\mu\) vor. Wir haben also zunächst in \(A\) den Complex \(a_\nu^{(\mu)}A_\nu^{(\mu)}\) und ferner die verschiedenen Complexe \(a_\nu^{(i)}a_k^{(\mu)}A_{\nu,k}^{i,\mu}\), worin \(i\) jeden von \(\mu\) verschiedenen und \(k\) jeden von \(\nu\) verschiedenen Index bedeuten kann.

\gesperrt{VI. Wir können daher setzen:}
\begin{equation}
        A=a_\nu^{(\mu)}A_\nu^{(\mu)}+
        \sum^{i,k} a_\nu^{(i)}a_k^{(\mu)}A_{\nu,k}^{i,\mu},
\tag{9}
\end{equation}
oder nach IV.
\begin{equation}
        A=a_\nu^{(\mu)}A_\nu^{(\mu)}-
        \sum^{i,k} a_\nu^{(i)}a_k^{(\mu)}A_{\nu,k}^{\mu,i}.
\tag{10}
\end{equation}

Wir bemerken zu diesem Satze noch, dass \(A_{\nu,k}^{\mu,i}\) die dem Elemente \(a_k^{(i)}\) entsprechende \gesperrt{erste Unterdeterminante} der \((n-1)\)reihigen Determinante \(A_\nu^{(\mu)}\) ist; denn \(A_{\nu,k}^{\mu,i}\) ist der Coëfficient
%CONT>>
